% Build with XeLaTeX.
\usepackage{iftex}
\ifXeTeX\else
  \PackageError{resume-template}{This template requires XeLaTeX}{Compile with xelatex or latexmk -xelatex.}
\fi

\usepackage{fontspec}
\usepackage{xeCJK}

% Font fallbacks keep the template portable across Overleaf, macOS and Windows.
\IfFontExistsTF{TeX Gyre Heros}{
  \setmainfont{TeX Gyre Heros}[UprightFont=*,BoldFont=* Bold,ItalicFont=* Italic,BoldItalicFont=* Bold Italic]
  \setsansfont{TeX Gyre Heros}[UprightFont=*,BoldFont=* Bold,ItalicFont=* Italic,BoldItalicFont=* Bold Italic]
}{
  \IfFontExistsTF{Helvetica}{
    \setmainfont{Helvetica}
    \setsansfont{Helvetica}
  }{
    \setmainfont{Latin Modern Sans}
    \setsansfont{Latin Modern Sans}
  }
}
\IfFontExistsTF{Latin Modern Mono}{
  \setmonofont{Latin Modern Mono}
}{
  \IfFontExistsTF{Menlo}{
    \setmonofont{Menlo}
  }{
    \setmonofont{Courier New}
  }
}

\IfFontExistsTF{Noto Serif CJK SC}{
  \setCJKmainfont{Noto Serif CJK SC}
}{
  \IfFontExistsTF{Source Han Serif SC}{
    \setCJKmainfont{Source Han Serif SC}
  }{
    \IfFontExistsTF{Songti SC}{
      \setCJKmainfont{Songti SC}
    }{
      \setCJKmainfont{SimSun}
    }
  }
}

\IfFontExistsTF{Noto Sans CJK SC}{
  \setCJKsansfont{Noto Sans CJK SC}
}{
  \IfFontExistsTF{Source Han Sans SC}{
    \setCJKsansfont{Source Han Sans SC}
  }{
    \IfFontExistsTF{PingFang SC}{
      \setCJKsansfont{PingFang SC}
    }{
      \setCJKsansfont{SimHei}
    }
  }
}

% Layout and basics.
\usepackage[
  ignoreheadfoot,
  top=1cm,
  bottom=1cm,
  left=1cm,
  right=1.25cm,
  footskip=1.25cm
]{geometry}
\usepackage{hyperref}
\usepackage[dvipsnames]{xcolor}

\definecolor{primaryColor}{RGB}{0,0,0}
\hypersetup{
  colorlinks=true,
  urlcolor=primaryColor,
  pdfborder={0 0 0}
}

\raggedright
\pagestyle{empty}
\setcounter{secnumdepth}{0}
\setlength{\parindent}{0pt}
\setlength{\topskip}{0pt}
\setlength{\columnsep}{0.15cm}
\pagenumbering{gobble}

% Reduce overfull boxes in mixed Chinese/English content.
\emergencystretch=2em
\XeTeXlinebreaklocale "zh"
\XeTeXlinebreakskip=0pt plus 1pt

% Section style without extra package dependencies.
\renewcommand{\section}[1]{%
  \par\vspace{0.20cm}%
  {\large\bfseries #1\par}%
  \vspace{1pt}\hrule\vspace{0.20cm}%
}

% Bullets and list spacing.
\renewcommand\labelitemi{\textbullet}
\newenvironment{highlights}{
  \begin{itemize}
    \setlength{\topsep}{0.10cm}
    \setlength{\parsep}{0.10cm}
    \setlength{\partopsep}{0pt}
    \setlength{\itemsep}{0pt}
}{\end{itemize}}

% Body block: keep a consistent right-side date/location gutter.
\newenvironment{onecolentry}{%
  \par\noindent
  \begin{list}{}{%
    \setlength{\leftmargin}{0pt}%
    \setlength{\rightmargin}{0.20\linewidth}%
    \setlength{\listparindent}{0pt}%
    \setlength{\itemindent}{0pt}%
    \setlength{\parsep}{0.2em}%
    \setlength{\topsep}{0pt}%
  }%
  \item\relax
}{%
  \end{list}%
  \par
}

% Entry heading: left content wraps, right side keeps date/location aligned.
\newcommand{\twocolentryright}{}
\newenvironment{twocolentry}[1]{%
  \def\twocolentryright{#1}%
  \par\noindent
  \begin{minipage}[t]{0.78\linewidth}\setlength{\parskip}{0.2em}
}{%
  \end{minipage}\hfill
  \begin{minipage}[t]{0.20\linewidth}\raggedleft
    \small \twocolentryright
  \end{minipage}%
  \par
}

\newcommand{\entrysep}{\vspace{0.20cm}}
